% GENERATED FILE. Edit canonical lessons, then run npm run generate:books.
% canonical-source-hash: fnv1a64:e5bbb300926118be
% canonical-lessons: ES-C733-gripe, ES-C733-resfriado, ES-C733-constipado, ES-C733-contagiar, ES-C733-estornudar

\chapter{Gripe y resfriado}
\label{ch:es-gripe-y-resfriado}
\hlchaptermodality{733}

\begin{chapteropening}
\textbf{By the end of this chapter:} \emph{I can say that I have a cold or the flu and that I keep sneezing.}

Complete the last lesson of chapter 733: I can say that I have a cold or the flu and that I keep sneezing.
\end{chapteropening}

\section[la gripe]{la gripe — the flu}

\label{lesson:ES-C733-gripe}



\begin{quote}
Before the new one: say the Spanish for to fall ill, then the Spanish for to get worse.
\end{quote}

\subsection*{You'll want to know: la gripe}
\textbf{la gripe} — \textquotedblleft{}the flu\textquotedblright{}.

\begin{quote}
Estoy en cama con gripe.
\end{quote}

\begin{quote}
\emph{I am in bed with the flu.}
\end{quote}

In Mexico and Colombia people say \emph{la gripa}.

\begin{cousinweb}
\emph{gripe} was borrowed from French \emph{grippe} (EtymDB), from \emph{gripper} \textquotedblleft{}to seize\textquotedblright{}, which comes from Frankish \emph{grīpan}: an illness that grips you.
\end{cousinweb}

\begin{grammarlens}[title={What you've built}]
One more word for telling, explaining and describing at B1.
\end{grammarlens}

\subsection*{Your turn}
\begin{itemize}
\raggedright
  \item \emph{Say it:} \emph{la gripe}
  \item \emph{Say it:} \emph{la gripe}, once more
  \item \emph{Say it:} a sentence of your own with \emph{la gripe}
  \item \emph{Recall:} say the Spanish for to fall ill, then the Spanish for to get worse, then say \emph{la gripe} again
\end{itemize}

\begin{tcolorbox}[breakable,colback=teal!4,colframe=teal!35!black,title={Before you move on}]
What does la gripe mean? (\textquotedblleft{}The flu\textquotedblright{}.) Say it once more.
\end{tcolorbox}
\section[el resfriado]{el resfriado — the cold (illness)}

\label{lesson:ES-C733-resfriado}



\begin{quote}
Before the new one: say the Spanish for to get worse, then the Spanish for the flu.
\end{quote}

\subsection*{You'll want to know: el resfriado}
\textbf{el resfriado} — \textquotedblleft{}the cold (illness)\textquotedblright{}.

\begin{quote}
Tengo un resfriado fuerte.
\end{quote}

\begin{quote}
\emph{I have a bad cold.}
\end{quote}

In Spain also \emph{el catarro}; in parts of South America \emph{el resfrío}.

\begin{cousinweb}
\emph{resfriado} is built on \emph{resfriarse} \textquotedblleft{}to catch cold\textquotedblright{}, from \emph{frío} \textquotedblleft{}cold\textquotedblright{}, which is usually traced to Latin \emph{frīgidus}.
\end{cousinweb}

\begin{grammarlens}[title={What you've built}]
One more word for telling, explaining and describing at B1.
\end{grammarlens}

\subsection*{Your turn}
\begin{itemize}
\raggedright
  \item \emph{Say it:} \emph{el resfriado}
  \item \emph{Say it:} \emph{el resfriado}, once more
  \item \emph{Say it:} a sentence of your own with \emph{el resfriado}
  \item \emph{Recall:} say the Spanish for to get worse, then the Spanish for the flu, then say \emph{el resfriado} again
\end{itemize}

\begin{tcolorbox}[breakable,colback=teal!4,colframe=teal!35!black,title={Before you move on}]
What does el resfriado mean? (\textquotedblleft{}The cold (illness)\textquotedblright{}.) Say it once more.
\end{tcolorbox}
\section[constipado]{constipado — having a cold, bunged up}

\label{lesson:ES-C733-constipado}



\begin{quote}
Before the new one: say the Spanish for the flu, then the Spanish for the cold.
\end{quote}

\subsection*{You'll want to know: constipado}
\textbf{constipado} — \textquotedblleft{}having a cold, bunged up\textquotedblright{}.

\begin{quote}
Estoy constipado y no puedo respirar bien.
\end{quote}

\begin{quote}
\emph{I have a cold and cannot breathe properly.}
\end{quote}

A false friend, mainly in Spain: \emph{estar constipado} is to have a cold. Constipated is \emph{estreñido}.

\begin{cousinweb}
\emph{constipado} comes from Latin \emph{cōnstīpātus} \textquotedblleft{}crowded or pressed closely together\textquotedblright{} (EtymDB), from \emph{stīpāre} \textquotedblleft{}to pack\textquotedblright{}: a head packed tight.
\end{cousinweb}

\begin{grammarlens}[title={What you've built}]
One more word for telling, explaining and describing at B1.
\end{grammarlens}

\subsection*{Your turn}
\begin{itemize}
\raggedright
  \item \emph{Say it:} \emph{constipado}
  \item \emph{Say it:} \emph{constipado}, once more
  \item \emph{Say it:} a sentence of your own with \emph{constipado}
  \item \emph{Recall:} say the Spanish for the flu, then the Spanish for the cold, then say \emph{constipado} again
\end{itemize}

\begin{tcolorbox}[breakable,colback=teal!4,colframe=teal!35!black,title={Before you move on}]
What does constipado mean? (\textquotedblleft{}Having a cold, bunged up\textquotedblright{}.) Say it once more.
\end{tcolorbox}
\section[contagiar]{contagiar — to infect, to pass on}

\label{lesson:ES-C733-contagiar}



\begin{quote}
Before the new one: say the Spanish for the cold, then the Spanish for having a cold.
\end{quote}

\subsection*{You'll want to know: contagiar}
\textbf{contagiar} — \textquotedblleft{}to infect, to pass on\textquotedblright{}.

\begin{quote}
No quiero contagiarte el resfriado.
\end{quote}

\begin{quote}
\emph{I do not want to give you my cold.}
\end{quote}

\emph{Contagiarse} — to catch something. It is also used for feelings: \emph{contagia su alegría}.

\begin{cousinweb}
\emph{contagiar} is built on \emph{contagio}, usually described as a learned borrowing from Latin \emph{contāgium} \textquotedblleft{}contact, infection\textquotedblright{}, from \emph{tangere} \textquotedblleft{}to touch\textquotedblright{}.
\end{cousinweb}

\begin{grammarlens}[title={What you've built}]
One more word for telling, explaining and describing at B1.
\end{grammarlens}

\subsection*{Your turn}
\begin{itemize}
\raggedright
  \item \emph{Say it:} \emph{contagiar}
  \item \emph{Say it:} \emph{contagiar}, once more
  \item \emph{Say it:} a sentence of your own with \emph{contagiar}
  \item \emph{Recall:} say the Spanish for the cold, then the Spanish for having a cold, then say \emph{contagiar} again
\end{itemize}

\begin{tcolorbox}[breakable,colback=teal!4,colframe=teal!35!black,title={Before you move on}]
What does contagiar mean? (\textquotedblleft{}To infect, to pass on\textquotedblright{}.) Say it once more.
\end{tcolorbox}
\section[estornudar]{estornudar — to sneeze}

\label{lesson:ES-C733-estornudar}



\begin{quote}
Before the new one: say the Spanish for having a cold, then the Spanish for to infect.
\end{quote}

\subsection*{You'll want to know: estornudar}
\textbf{estornudar} — \textquotedblleft{}to sneeze\textquotedblright{}.

\begin{quote}
El polvo me hace estornudar.
\end{quote}

\begin{quote}
\emph{Dust makes me sneeze.}
\end{quote}

People say \emph{¡Salud!} or \emph{¡Jesús!} when someone sneezes.

\begin{cousinweb}
\emph{estornudar} is inherited from Latin \emph{sternūtāre} \textquotedblleft{}to sneeze\textquotedblright{} (EtymDB).
\end{cousinweb}

\begin{grammarlens}[title={What you've built}]
One more word for telling, explaining and describing at B1.
\end{grammarlens}

\subsection*{Your turn}
\begin{itemize}
\raggedright
  \item \emph{Say it:} \emph{estornudar}
  \item \emph{Say it:} \emph{estornudar}, once more
  \item \emph{Say it:} a sentence of your own with \emph{estornudar}
  \item \emph{Recall:} say the Spanish for having a cold, then the Spanish for to infect, then say \emph{estornudar} again
\end{itemize}

\begin{tcolorbox}[breakable,colback=teal!4,colframe=teal!35!black,title={Before you move on}]
What does estornudar mean? (\textquotedblleft{}To sneeze\textquotedblright{}.) Say it once more.
\end{tcolorbox}
